% Generated draft: requires mathematical and source review.
\documentclass[11pt]{article}
\usepackage[margin=25mm]{geometry}
\usepackage{fontspec}
\setmainfont{Latin Modern Roman}
\usepackage{amsmath,amssymb,mathtools,unicode-math}
\setmathfont{Latin Modern Math}
\usepackage{enumitem,needspace,xurl,hyperref}
\setlength{\parindent}{0pt}
\setlength{\parskip}{5pt}
\setlength{\emergencystretch}{3em}
\newcommand{\N}{\mathbb{N}}
\newcommand{\Z}{\mathbb{Z}}
\newcommand{\Q}{\mathbb{Q}}
\newcommand{\R}{\mathbb{R}}
\newcommand{\C}{\mathbb{C}}
\newcommand{\F}{\mathbb{F}}
\newcommand{\A}{\mathbb{A}}
\newcommand{\PP}{\mathbb{P}}
\newcommand{\E}{\mathbb{E}}
\newcommand{\eps}{\varepsilon}
\newcommand{\dd}{\,\mathrm{d}}
\begin{document}
\section*{UnsolvedMath ID: 20003110}
\textbf{Source problem code:} catalogue entry
\subsection{Definitions and mathematical statement}
Let $S$ be a finite-type bordered hyperbolic surface with nonempty geodesic boundary, and let $\mathcal T(S)$ denote the reduced Teichm\"uller space used in the source. Boundary lengths may vary and boundary parametrizations are not retained. In the usual arc-metric convention, let $\mathcal A(S)$ be essential properly embedded arc classes with endpoints allowed to slide on the boundary, and include essential closed curves and boundary curves in the test set $\mathcal L(S)$. For each marked structure, arcs are represented by perpendicular geodesic arcs. Set
\[
d_A(X,Y)=\log\sup_{\alpha\in\mathcal L(S)}\frac{\ell_Y(\alpha)}{\ell_X(\alpha)}.
\]
The exact arc/curve and marking conventions require agreement with the cited reduced-space construction.

The original problem has two parts: is this asymmetric arc metric a length metric, and can every ordered pair be joined by a Thurston geodesic in the double that stays inside the symmetric doubled locus? More precisely, let $S^d$ and $X^d$ be the topological and hyperbolic doubles across the boundary, and let
\[
j:\mathcal T(S)\longrightarrow\mathcal T(S^d),\qquad j(X)=X^d.
\]
The source's notation $\Gamma\subset\mathcal T(S)$ inside $\mathcal T(S^d)$ means $\Gamma\subset j(\mathcal T(S))$.

The aid attributes to Alessandrini--Disarlo the following affirmative result: for every ordered pair $X,Y$, there is a forward unit-speed arc geodesic
\[
\gamma:[0,D]\to\mathcal T(S),\qquad D=d_A(X,Y),\quad
\gamma(0)=X,\ \gamma(D)=Y,
\]
\[
d_A(\gamma(s),\gamma(t))=t-s\qquad(0\leq s\leq t\leq D),
\]
which is a finite concatenation of generalized stretch segments. Thus the forward distance is realized by a path and equals its intrinsic forward length. This does not claim uniqueness or that one ordinary stretch segment suffices.

The aid also invokes the Liu--Papadopoulos--Su--Th\'eret doubling identity
\[
d_{\mathrm{Th}}(X^d,Y^d)=d_A(X,Y).
\]
Consequently $j\circ\gamma$ is a forward unit-speed Thurston geodesic with its entire image in $j(\mathcal T(S))$, answering the second question as well. In fact, for each fixed ordered pair, composition with $j$ gives a bijection between all parameterized forward unit-speed arc geodesics from $X$ to $Y$ and all parameterized forward unit-speed Thurston geodesics from $X^d$ to $Y^d$ whose images lie in this doubled locus. Its inverse is pointwise restriction by $j^{-1}$.

Thus the ambient-geodesic clause is exactly arc-metric geodesicity under the doubling identity. It is stronger than merely asserting a length-space axiom in the absence of a separate minimizing-path existence argument; length alone does not assert that the infimum is attained. The reported generalized-stretch theorem supplies that attainment in the stated reduced bordered setting.
\subsection{Short English statement}
The supplied literature answers both questions affirmatively: the arc metric has minimizing geodesics, and doubling identifies exactly those geodesics with Thurston geodesics that remain in the symmetric doubled locus.
\subsection{Sources}
\begin{itemize}
\item Local UnsolvedMath registry, permanent ID $20003110$, source code \texttt{catalogue entry}; canonical \texttt{aim-topology-notes.json}, \texttt{notes[197]}.
\item Weixu Su's 2014 AIM list \emph{Problems on Thurston Metric}, Problem 3.7, workshop \emph{Lipschitz metric on Teichmueller space}: \url{https://aimath.org/pastworkshops/lipschitzteichproblems.pdf}. The supplied verification says preceding text fixes the notation; original remarks and literature field are empty. Trailing dots are extraction punctuation.
\item Alessandrini--Disarlo for forward arc geodesics as finite concatenations of generalized stretch segments; Liu--Papadopoulos--Su--Th\'eret for the doubling distance identity, as named in the aid. No titles, theorem numbers or authoritative full bibliographic entries are supplied.
\item Research attempt 1: \texttt{full\_solution}, status \texttt{solved}, classification \texttt{SOLVED-IN-LITERATURE}, rigor \texttt{proved}, novelty confidence low. Its geodesic-bijection corollary and literature assessment are machine-generated aid, not peer reviewed, and require independent verification before citation.
\end{itemize}
\end{document}
